\section{A local comparison}\label{sec:local}

The varieties in this section and the next are over \(\C\).
We use the local multiplicity comparison underlying the
separated-weight method of \cite{OAIpi}; see also
\cite[Lemma 3]{LJZ2026}.

A normal basis to a smooth subvariety means a collection of ambient
vectors whose images form a basis modulo its tangent space. For an
irreducible subvariety, all statements about these bases are at a
general smooth point.

\begin{lemma}\label{lem:local}
Let affine coordinates have positive rational degree weights
\(W_a\). Let commuting analytic vector fields \(D_b\) form a frame
on an open set and have positive costs \(V_b\). Suppose polynomials
\(f_\ell\) have weighted degree at most \(N>0\), and \(S\) is a
codimension-\(k\) irreducible component, with \(k\geq1\), of their common zero set in
that open set. If all further frame derivatives of every \(f_\ell\)
of cost less than \(\epsilon N\) vanish on \(S\), then for a
coordinate normal basis \(A\) and a frame normal basis \(B\),
\begin{equation}\label{eq:comparison}
 \prod_{a\in A}W_a
 \leq k!\epsilon^{-k}\prod_{b\in B}V_b
 \leq k^k\epsilon^{-k}\prod_{b\in B}V_b.
\end{equation}
\end{lemma}
\begin{proof}
Choose a smooth point of \(S\) with both basis properties and away
from every other component. The commuting flows indexed by \(B\),
started on \(S\), give analytic coordinates \((z,b)\). Here \(z\)
is on \(S\) and \(b\) is transverse. The derivative hypothesis says
that each equation has an expansion
\[
 f_\ell(z,b)=\sum_{\sum_{j\in B}V_j\alpha_j\geq\epsilon N}
                         f_{\ell,\alpha}(z)b^\alpha.
\]
On the affine coordinate slice obtained by fixing coordinates outside
\(A\), the map to \(b\) is locally invertible. Substitution of its
analytic functions \(z=z(b)\) cannot decrease the weight of a surviving
monomial. Consequently, the quotient by the restricted equations has
length at least the number of nonnegative integer \(\alpha\) with
\(\sum_{j\in B}V_j\alpha_j<\epsilon N\). The unit cubes with those
lower corners cover the corresponding open real simplex, so this number
is at least
\[
                 \frac{(\epsilon N)^k}{k!\prod_{j\in B}V_j}.
\]

The common zero on the coordinate slice is isolated. Generic \(k\)
constant linear combinations of the restricted equations still have
an isolated common zero there; their quotient length bounds the previous
length from above. Weighted B\'ezout gives the upper bound
\(N^k/\prod_{a\in A}W_a\). One can obtain precisely this bound from
ordinary isolated-zero B\'ezout: choose an integer \(M_0\) with
\(M_0W_a\) positive integers, and substitute
\(x_a=\lambda_a+z_a^{M_0W_a}\), choosing \(\lambda_a\) away from the
coordinates of the point. The \(\prod M_0W_a\) preimages are unramified,
preserve local length, and the transformed equations have total degree
at most \(M_0N\). Comparison of the two length bounds proves the claim.
\end{proof}

\section{Two-base interpolation}\label{sec:geometry}

Fix positive integers \(m,K\), positive rational \(w_0,v_0,v,\sigma\) with
\(v<1\), and positive rational \(w_1,\ldots,w_m\). Index the two
base coordinates first and the auxiliary coordinates in increasing order.
Set
\[
 W=(w_0,w_0,w_1,\ldots,w_m),\qquad
 V=(v_0,v_0,w_1/v,\ldots,w_m/v),\qquad d=m+2.
\]
For an auxiliary multi-index \(\alpha\), write
\(w\alpha=\sum_i w_i\alpha_i\). Let \(\PW(H)\) be spanned by
\(Y^hZ^kX^\alpha\) with \(w_0(h+k)+w\alpha\leq H\).
Let \(\JV(H)\) retain the formal coefficients
\(t^{\ell_1}s^{\ell_2}u^\beta\) for which
\(v_0(\ell_1+\ell_2)+w\beta/v<H\).

Put \(\epsilon=\sigma/(m+3)\). Impose the separation inequalities
\begin{equation}\label{eq:separation}
 \prod_{a\in A}W_a>
       k^k\epsilon^{-k}\prod_{b\in B}V_b
\end{equation}
for every two index sets of the same cardinality
\(1\leq k\leq m+1\) whose largest differing auxiliary index lies
in \(A\). There are successive lower thresholds on the \(w_i\)
ensuring all these inequalities. Indeed, at the largest differing
auxiliary index the ratio contains the new weight linearly. Higher
common indices cancel except for fixed factors \(W_j/V_j=v\);
the other weights are already fixed. These thresholds do not involve
any rational slope.

For an explicit simultaneous choice put
\[
 L=\min(w_0,1),\quad R_0=\max(v_0,1/v),\quad
 D=1+\sum_{k=0}^{d}k^k\epsilon^{-k}(R_0/L)^k,
\]
using \(0^0=1\). Write \(x_0=x_1=1\) and \(x_{i+1}=w_i\)
for \(1\le i\le m\). It suffices that all \(x_i\ge1\) and
\(x_i>D\prod_{j<i}x_j\) for every auxiliary index \(i\ge2\).
Indeed \(W_j\ge Lx_j\) and \(V_j\le R_0x_j\). Cancel common
indices from the two products. At the largest differing index, the
growth condition exceeds the product of all earlier indices by the
factor \(D\); the remaining factors on its side are at least one.
The definition of \(D\) then gives \eqref{eq:separation} for every
cardinality simultaneously.

\begin{proposition}\label{prop:interpolation}
Suppose \eqref{eq:separation} and
\begin{align}
 (1+3\sigma)^d K(w_0/v_0)^2v^m&<1,\label{eq:volume}\\
 (1+\sigma)v&<1,\label{eq:shrink}\\
 w_0/v_0&>2(m+3)/\sigma.\label{eq:basegap}
\end{align}
Let \(r_i=p_i/q_i\) be reduced positive fractions with
\begin{equation}\label{eq:qcondition}
                         q_i v_0/w_0>1+\sigma.
\end{equation}
At the centers \((2^j,3^j,0,\ldots,0)\), \(0\leq j<K\), use
the substitutions
\begin{equation}\label{eq:jets}
 Y=2^j(1+t),\quad Z=3^j(1+s),\quad
 X_i=\log(1+s)-r_i\log(1+t)+u_i.
\end{equation}
Then \(\PW(H)\longrightarrow\bigoplus_{j<K}\JV(H)\) is
surjective for every sufficiently large \(H\) divisible by a fixed
positive integer. Its degree threshold may depend on all the fixed data.
\end{proposition}

\subsection{A curve inequality}

On the complete smooth model of an irreducible affine curve define
\[
 D_W(C)=\sum_P\max\left(0,-\frac{\ord_PY}{w_0},
        -\frac{\ord_PZ}{w_0},
        -\frac{\ord_PX_1}{w_1},\ldots,-\frac{\ord_PX_m}{w_m}\right).
\]
An identically zero coordinate contributes no pole. At a branch through
center \(j\), put \(t=Y/2^j-1\), \(s=Z/3^j-1\), and
\(u_i=X_i-\log(1+s)+r_i\log(1+t)\), and set
\[
 h_P=\min\left(\frac{\ord_Pt}{v_0},\frac{\ord_Ps}{v_0},
                         \frac{\ord_Pu_i}{w_i/v}:1\leq i\leq m\right).
\]
Orders of identically zero series are infinity. The formal coordinate
change is invertible, so \(h_P\) is positive and finite.
We prove
\begin{equation}\label{eq:curve}
                D_W(C)\geq(1+\sigma)\sum_{P\text{ centered}}h_P.
\end{equation}

Suppose it fails. In particular the curve meets a center and neither
\(Y\) nor \(Z\) vanishes identically. The leading dimensions of
polynomials of degree at most \(N\) and of the conditions requiring
vanishing below \((1+3\sigma)N\) are respectively
\[
 \frac{N^d}{d!\prod_a W_a},\qquad
 \frac{K(1+3\sigma)^dN^d}{d!\prod_a V_a}.
\]
Thus \eqref{eq:volume} supplies a nonzero such polynomial \(F_N\)
for all sufficiently large \(N\).

On \(YZ\ne0\) take the commuting frame
\[
 D_Y=Y\partial_Y-\sum_i r_i\partial_{X_i},\quad
 D_Z=Z\partial_Z+\sum_i\partial_{X_i},\quad D_i=\partial_{X_i}.
\]
It preserves the degree bound and becomes
\((1+t)\partial_t,(1+s)\partial_s,\partial_{u_i}\) in the jets.
A frame word of cost at most \(\sigma N\) applied to \(F_N\)
has branch order at least \((1+2\sigma)Nh_P\). Its total poles on
the curve are at most \(ND_W(C)\). Failure of \eqref{eq:curve}
would make the total zeros larger, unless that derivative vanished
identically on the curve. Therefore every such word does vanish.

Consider the zero sets of frame derivatives of costs at most
\(r\epsilon N\), \(0\leq r\leq m+3\). Starting from the last,
choose nested irreducible components containing the curve. Their
dimensions lie between one and \(m+1\), so two adjacent ones coincide.
Call the persistent component \(S\). Further derivatives of the earlier
equations of cost less than \(\epsilon N\) vanish on \(S\).
Lemma~\ref{lem:local} applies to every pair of its normal bases.

If all auxiliary coordinate directions are tangent to \(S\), its normal
space has dimension \(k=1\) or \(2\). Coordinate and frame normal bases
then use only base indices. The comparison gives
\(w_0/v_0\leq k/\epsilon\leq2/\epsilon\), contrary to
\eqref{eq:basegap}.

Otherwise choose the largest nontangent auxiliary direction \(i\).
Some coordinate normal basis contains it. Higher auxiliary directions
are zero in the normal quotient for both frames. Hence
\eqref{eq:separation} forces every frame normal basis to contain \(i\).
The other frame vectors span the hyperplane
\(\ker(dX_i+r_i\,dY/Y-dZ/Z)\); their images fail to span the normal
quotient. The tangent space of \(S\) is therefore in this hyperplane.
On its function field we have
\begin{equation}\label{eq:rigidity}
             d(q_iX_i)=d\log(Z^{q_i}/Y^{p_i}).
\end{equation}
Both rational functions in this identity are constant. To see this,
if the character were nonconstant, choose a smooth point where it and
\(X_i\) are regular and its differential is nonzero. Intersect \(S\)
with affine hyperplanes through that point whose common tangent line
is not in the kernel of this differential. The curve component through
the point is smooth there, and the character remains nonconstant on it.
Restrict \eqref{eq:rigidity} to this curve and take a complete smooth model.
At one of its zeros or poles its logarithmic differential has nonzero
residue, whereas an exact differential has residue zero. Once the
character is constant, characteristic zero makes \(X_i\) constant too.
Clearing denominators transfers these polynomial relations to the
original curve. The center character values
\((3^{q_i}/2^{p_i})^j\) are pairwise different, by unique factorization.
The curve therefore meets at most one center.

If \(Y\) is nonconstant on that curve, choose integers \(a,b\) with
\(a q_i+b p_i=1\). Writing \(c=Z^{q_i}/Y^{p_i}\ne0\) and
\(T=Y^aZ^b\), we have \(Y=c^{-b}T^{q_i}\). At every branch through
the possible center, \(T\) has the same nonzero residue \(T_j\).
The difference-of-powers factor has nonzero residue
\(q_iT_j^{q_i-1}\), whence
\[
 \ord_P(T-T_j)=\ord_P(Y-2^j)\geq v_0h_P.
\]
The pole degree of \(Y\) is \(q_i\deg(T)\). Zero--pole equality now
implies
\[
 D_W(C)\geq\frac{q_i}{w_0}\deg(T)
       \geq\frac{q_i v_0}{w_0}\sum_Ph_P
       >(1+\sigma)\sum_Ph_P,
\]
a contradiction. If \(Y\) is constant, \(Z^{q_i}=cY^{p_i}\) makes
\(Z\) constant. Some \(X_l\) is nonconstant, and the two logarithmic
corrections vanish along the curve. Therefore
\[
 (w_l/v)\sum_Ph_P\leq\deg(X_l)\leq w_lD_W(C),
\]
contradicting \eqref{eq:shrink}. Curves missing all centers have contact
sum zero. This proves \eqref{eq:curve} in every case.

\subsection{Obtaining bounded-degree polynomials}

Choose an integer \(R\) such that all \(R/W_a,R/V_a\) are positive
integers and \(R\geq\max_a W_a\). Embed the affine space by its monomials of weighted degree at
most \(R\), and let \(X\) be the projective closure with
\(L=\OO_X(1)\). The constant-coordinate chart is the original affine
space. Pure powers and the constant monomial show on any normalized
curve meeting this chart that
\begin{equation}\label{eq:degree}
                          \deg L=R D_W(C).
\end{equation}

Truncate both logarithms in auxiliary coordinate \(i\) with first
omitted degree \(T_i\) such that \(T_iv_0>w_i/v\), obtaining regular
coordinates \(t,s,u'_i\). Every omitted term has order greater than
\((w_i/v)h_P\), so the minimum defining \(h_P\) is unchanged. The
coordinate change and its inverse preserve the entire weighted filtration.
At each center form
\[
 I_j=(t^{R/v_0},s^{R/v_0},(u'_1)^{Rv/w_1},\ldots,(u'_m)^{Rv/w_m}).
\]
Patch with the unit ideal elsewhere to get a coherent ideal \(I\).
Its branch order is \(Rh_P\).

Let \(p:X'\to X\) be the ordinary blow-up of \(I\), and write
\(I\OO_{X'}=\OO_{X'}(-E)\), \(A=p^*L\). On a noncontracted curve
whose image meets the chart, \eqref{eq:curve} and \eqref{eq:degree}
give nonnegative degree for \(A-(1+\sigma)E\). Curves outside the
chart miss \(E\), and on contracted curves \(-E\) has positive degree
by relative ampleness. Thus \(A-(1+\sigma)E\) is nef.
For some integer \(a>1\), \(aA-E\) is ample, and
\[
 A-E=\frac{a-1}{a(1+\sigma)-1}\bigl(A-(1+\sigma)E\bigr)
       +\frac{\sigma}{a(1+\sigma)-1}(aA-E)
\]
is ample as well. These facts apply to the projective schemes in question;
smoothness of \(X'\) is not needed; see \cite[Corollary 1.4.10]{Lazarsfeld}
and \cite{StacksBlowup}.

For large \(n\), the graded-piece comparison and relative vanishing
for the Rees algebra yield
\(p_*\OO_{X'}(-nE)=I^n\) and
\(R^qp_*\OO_{X'}(-nE)=0\) for \(q>0\)
\cite{StacksGraded,StacksComparison}. Projection, Leray and Serre
vanishing give \(H^1(X,I^n\otimes L^n)=0\). Sections of \(L^n\)
therefore realize all local classes modulo \(I^n\).
Each generator has formal weighted order at least \(R\), so \(I^n\)
is contained in the ideal of weight at least \(nR\); the desired packet
quotient is smaller. Completion does not change these finite-length
quotients. Finally, Serre vanishing for the ideal sheaf of the projective
embedding lifts every section of \(L^n\), for large \(n\), to a
degree-\(n\) homogeneous polynomial in its monomial coordinates. On the
affine chart this has weighted degree at most \(nR\).
Taking \(H=nR\) proves Proposition~\ref{prop:interpolation}.
